\section{Design Principles of CARP}

We have designed a streaming in-situ partitioner for large parallel applications, CARP, which intercepts writes from applications, transforms them, and persists them in a range query-friendly layout. A logical view of CARP's partitioning process is shown in \cref{fig:carp:des-layout}. CARP has no impact on application runtime, as it uses spare CPU and network capacity to transform data at a rate necessary to saturate storage. This partitioned output requires no post-processing and can be queried directly once the application terminates.

\figdeslayout

CARP augments conventional range partitioning with novel techniques to achieve its write and query performance goals. We build upon a simple insight --- data does not need to be fully sorted to achieve the performance benefits of a sorted, clustered index. It is sufficient to partially order data in a way that allows queries to be processed via a small number of large reads. CARP exploits this relaxation to achieve its write performance and adaptivity goals.

We employ a greedy approach to generate best-effort range partitions in a single streaming pass so as to not require expensive reorganization of written data. To generate balanced partitions, CARP employs a novel primitive to construct and monitor the global key distribution. This distribution is divided into partitions with equal areas under the curve, and each partition is assigned to a rank. In the common case, all ranks shuffle data to corresponding partitions. As the workload's key distribution drifts, CARP recomputes partition boundaries and continues shuffling. Shuffled data is written to disk via a storage backend called \emph{KoiDB} --- KoiDB further optimizes CARP output to improve range query performance. Next, we describe how CARP achieves the design goals laid out in the previous section.

\textbf{Write Performance}. To ingest data at storage layer throughput, CARP is designed to not divert any storage bandwidth for reorganization.
Even when partitions need to be recomputed, CARP does not touch previously written data --- it simply records a change in partition boundaries on each rank and continues appending data. CARP also employs standard storage optimization practices such as batching small writes into large immutable units called \emph{SSTables} and having its on-disk layout be an append-only log.

\textbf{Query Performance}. Efficient queries in CARP are enabled by 1) creating a highly partitioned layout, and 2) enabling efficient retrieval via large sequential writes. CARP creates one partition per application rank, automatically ensuring more data partitions as application scale increases. A selective range query only needs to retrieve a fraction of the total data from overlapping partitions. In addition, data is written in large SSTables, which can be read efficiently via large read requests. CARP merges its partially ordered SSTs at query time via a merge-sort operation, which is cheap compared to the I/O cost of retrieving data.

\textbf{Adaptivity}. To adapt to skewed and dynamic key distributions, we recompute CARP partitions as the incoming distribution changes. Unlike conventional databases, we do not repartition data that has already been written to disk. We simply record the change in partition boundaries and continue appending new data. Over the lifetime of a run, the range allocated to a CARP partition changes multiple times. As a corollary, data belonging to one point in the range may end up in different partitions at different times in the run. Instead of incurring write amplification to consolidate this data, we simply postpone this consolidation (also called \emph{compaction}) to query time.
